% Source: https://artofproblemsolving.com/community/c3h2787023_question
% Stable identity: aops-topic-2787023-q01
(a) Write $P=(x_0,y_0)$ and $\mathbf n=(a,b)$, a normal vector of the line. The image $P'$ satisfies two conditions: $PP'$ is perpendicular to the line, so $P'=P+t\mathbf n$ for some $t$; and the midpoint $P+\dfrac t2\mathbf n$ lies on the line:
\[a\left(x_0+\dfrac{ta}2\right)+b\left(y_0+\dfrac{tb}2\right)+c=0\quad\Longrightarrow\quad t=-\dfrac{2(ax_0+by_0+c)}{a^2+b^2},\]
which gives the stated coordinates.

(b) Translate so that the centre is the origin: $(u,v)=(x_0-h,y_0-k)=(r\cos\varphi,r\sin\varphi)$. Rotating gives $(r\cos(\varphi+\theta),r\sin(\varphi+\theta))$, which by the addition formulas is $(u\cos\theta-v\sin\theta,\ u\sin\theta+v\cos\theta)$. Translating back adds $(h,k)$.

(c) The line $y=2x$ is $2x-y=0$, so $a=2$, $b=-1$, $c=0$ and $ax_0+by_0+c=10$. The reflection is
\[\left(5-\dfrac{2\cdot2\cdot10}{5},\ 0+\dfrac{2\cdot10}{5}\right)=\boxed{(-3,4)}.\]
(Check: the midpoint $(1,2)$ is on the line, and $(-8,4)$ is perpendicular to the direction $(1,2)$.)

For the rotation, $(u,v)=(3,-1)$ and $\theta=90^\circ$, so $(u\cos\theta-v\sin\theta,\ u\sin\theta+v\cos\theta)=(1,3)$, and the image is $(2,1)+(1,3)=\boxed{(3,4)}$.
\begin{asy}
usepackage("amsmath");
size(6.5cm);
defaultpen(fontsize(10pt));
draw((-4,0)--(6.3,0),Arrow(4));draw((0,-1.6)--(0,6.6),Arrow(4));
draw((-0.8,-1.6)--(3.2,6.4));label("$y=2x$",(3.2,6.4),E);
pair P=(5,0),P1=(-3,4),Q=(2,1),P2=(3,4),Mi=(1,2);
draw(P--P1,dashed);draw(Q--P,dotted);draw(Q--P2,dotted);
draw(arc(Q,1.1,-18.43,71.57),Arrow(4));
dot(P);dot(P1);dot(Q);dot(P2);dot(Mi);
label("$P(5,0)$",P,S);label("$P'(-3,4)$",P1,N);label("$(1,2)$",Mi,W);label("$Q(2,1)$",Q,S);label("$P''(3,4)$",P2,E);label("$90^\circ$",Q+1.4*dir(30),E);
\end{asy}
